\documentclass[../../main.tex]{subfiles}

\begin{document}

\chapter[Parser Generator dan Semantic Actions]{Parser Generator dan Semantic Actions}
\label{chap:parser-generator}

\begin{outcomebox}
\textbf{Sub-CPMK:} \subcpmk{4.3}, \subcpmk{4.4}\\[3pt]
Setelah menyelesaikan bab ini, mahasiswa mampu:
\begin{enumerate}[label=\arabic*.,leftmargin=*]
  \item menjelaskan arsitektur parser generator;
  \item mengimplementasikan parser subset bahasa C dengan Flex dan
        Bison beserta aksi semantik;
  \item membangun AST dari semantic action;
  \item mengintegrasikan Flex dan Bison dalam satu build;
  \item menguji parser dan AST dengan keluaran program serta pesan
        kesalahan sintaksis.
\end{enumerate}
\end{outcomebox}

\section*{Materi Pokok}
\begin{itemize}[leftmargin=*]
  \item Arsitektur parser generator.
  \item Deklarasi, aturan, dan konflik Bison.
  \item Semantic action dan konstruksi AST.
  \item Integrasi Flex--Bison dan build.
  \item Pengujian parser, AST, dan pesan kesalahan sintaksis.
\end{itemize}

\subfile{section-01}
\subfile{section-02}
\subfile{section-03}
\subfile{section-04}
\subfile{section-05}

\section*{Contoh}
\addcontentsline{toc}{section}{Contoh}

\begin{contoh}[Aturan ekspresi Bison]
Aturan
\begin{lstlisting}[language=Bison]
expr : expr T_PLUS expr
     | expr T_STAR expr
     | T_INT_LIT
     ;
\end{lstlisting}
dapat menimbulkan konflik shift-reduce karena precedence tidak
dideklarasikan. Dengan \code{\%left T\_PLUS} dan
\code{\%left T\_STAR}, Bison menyelesaikan konflik sesuai prioritas
perkalian.
\end{contoh}

\begin{contoh}[Membangun simpul AST]
Aksi
\begin{lstlisting}[language=Bison]
expr : expr T_PLUS expr
     {
         $$ = ast_new(AST_BINARY, $2.token);
         $$->left = $1;
         $$->right = $3;
     }
     ;
\end{lstlisting}
membuat simpul \code{BINARY} dengan token operator dari \code{\$2},
anak kiri dari \code{\$1}, dan anak kanan dari \code{\$3}.
\end{contoh}

\begin{contoh}[AST untuk \(a + b * c\)]
Aksi \code{term} membangun simpul \code{*} untuk \(b * c\). Aksi
\code{additive} membangun simpul \code{+} dengan anak \code{a} dan
simpul \code{*}. Hasilnya adalah AST pada
\figref{fig:ch13-ast}.
\end{contoh}

\begin{contoh}[Precedence dan associativity]
Deklarasi
\begin{lstlisting}[language=Bison]
%left T_PLUS T_MINUS
%left T_STAR T_SLASH T_PERCENT
\end{lstlisting}
membuat \code{*} lebih tinggi daripada \code{+}. Keduanya
kiri-asosiatif, sehingga \(a - b - c\) diurai sebagai
\((a - b) - c\). Untuk operator kanan-asosiatif, gunakan
\code{\%right}.
\end{contoh}

\begin{contoh}[Location tracking]
Dengan \code{\%locations}, Bison menyediakan \code{@1}, \code{@2},
dan \code{@\$} untuk lokasi simbol. Aksi dapat menyalin
\code{@\$.first\_line} ke simpul AST, sehingga diagnostik semantik
dapat menyebut baris yang tepat.
\end{contoh}

\section*{Aktivitas Pembelajaran dan Praktikum}
\addcontentsline{toc}{section}{Aktivitas Pembelajaran dan Praktikum}

\begin{table}[H]
\centering
\caption{Rencana aktivitas Bab XIII}
\label{tab:ch13-aktivitas}
\begin{tabularx}{\textwidth}{@{}llX@{}}
\toprule
\textbf{Sesi} & \textbf{Bentuk} & \textbf{Kegiatan} \\
\midrule
1 & Ceramah & Arsitektur parser generator dan Bison \\
2 & Live coding & Menulis grammar dan semantic action \\
3 & Latihan & Membaca laporan konflik Bison \\
4 & Praktikum & M3: Bison + AST + integrasi \\
\bottomrule
\end{tabularx}
\end{table}

\subsection*{M3 / Praktikum 3: Bison dan AST}
\textbf{Tujuan.} Membangun parser Bison untuk Subset-C dan
menghasilkan AST lengkap.

\textbf{Keluaran.}
\begin{enumerate}[label=\arabic*.,leftmargin=*]
  \item Grammar Bison untuk ekspresi, pernyataan, dan deklarasi.
  \item Semantic action yang membangun AST.
  \item Integrasi Flex--Bison dan driver.
  \item Minimal 20 kasus uji: valid dan tidak valid.
  \item Perbandingan dengan recursive descent.
\end{enumerate}

\textbf{Langkah kerja.}
\begin{enumerate}[label=\arabic*.,leftmargin=*]
  \item Lengkapi \file{src/parser.y}.
  \item Definisikan \code{\%union}, \code{\%token}, dan \code{\%type}.
  \item Tulis semantic action untuk setiap aturan.
  \item Integrasikan \code{yylex} dari lexer.
  \item Jalankan Bison dan perbaiki konflik.
  \item Uji dan bandingkan hasilnya dengan parser tangan.
\end{enumerate}

\textbf{Rubrik.}
\begin{table}[H]
\centering
\caption{Rubrik Praktikum 3}
\label{tab:ch13-rubrik}
\begin{tabularx}{\textwidth}{@{}lXr@{}}
\toprule
\textbf{Aspek} & \textbf{Indikator} & \textbf{Bobot} \\
\midrule
Kebenaran & Grammar dan AST benar & 40\% \\
Kualitas kode & Aksi ringkas, penamaan, komentar & 30\% \\
Kelengkapan & Ekspresi, pernyataan, deklarasi & 20\% \\
Pengujian & Kasus normal, batas, kesalahan & 10\% \\
\bottomrule
\end{tabularx}
\end{table}

\textbf{Definisi selesai.}
\begin{checklist}
  \item Grammar Bison tanpa konflik tak terduga.
  \item AST terbentuk untuk program valid.
  \item 20 kasus uji tersedia.
  \item Perbandingan dengan recursive descent terdokumentasi.
\end{checklist}

\section*{Latihan dan Refleksi}
\addcontentsline{toc}{section}{Latihan dan Refleksi}

\subsection*{Latihan}
\begin{enumerate}[label=\arabic*.,leftmargin=*]
  \item Jelaskan peran \code{\%token}, \code{\%type}, dan
        \code{\%union}.
  \item Tuliskan aturan Bison untuk ekspresi penjumlahan dan
        perkalian.
  \item Jelaskan arti \code{\$\$} dan \code{\$1}.
  \item Bagaimana semantic action membangun AST?
  \item Mengapa precedence dapat menyelesaikan konflik?
  \item Jelaskan integrasi \code{yylex} dan \code{yyparse}.
  \item Bandingkan Bison dan parser tangan dari sisi kontrol,
        produktivitas, dan diagnostik.
  \item Kapan parser tangan lebih tepat daripada parser generator?
  \item Tuliskan pesan kesalahan sintaksis untuk tiga kasus berbeda;
        setiap pesan harus memuat lokasi, penyebab, dan harapan token.
\end{enumerate}

\subsection*{Refleksi}
\begin{reflectionbox}
\begin{enumerate}[label=\arabic*.,leftmargin=*]
  \item Apakah Anda dapat membaca laporan konflik Bison?
  \item Bagian mana dari semantic action yang paling menantang?
  \item Alat mana yang akan Anda pilih untuk proyek nyata? Mengapa?
\end{enumerate}
\end{reflectionbox}

\section*{Asesmen}
\addcontentsline{toc}{section}{Asesmen}

Bab XIII memuat Praktikum 3 dengan bobot 7\% sesuai RPS. Asesmen
mencakup grammar, AST, integrasi, dan pengujian.

\begin{table}[H]
\centering
\caption{Instrumen asesmen Bab XIII}
\label{tab:ch13-asesmen}
\begin{tabularx}{\textwidth}{@{}lXl@{}}
\toprule
\textbf{Instrumen} & \textbf{Indikator} & \textbf{Bobot} \\
\midrule
Praktikum 3 & Grammar Bison, AST, integrasi, tes & 7\% \\
Kuis & Semantic action dan konflik & formatif \\
Review kode & Kualitas grammar dan aksi & bagian P3 \\
\bottomrule
\end{tabularx}
\end{table}

Kriteria kelulusan: grammar Bison membangun AST yang benar, konflik
ditangani secara eksplisit, dan integrasi dengan lexer berjalan.

\section*{Checklist Pencapaian Kompetensi}
\addcontentsline{toc}{section}{Checklist Pencapaian Kompetensi}

\begin{checklist}
  \item Saya dapat menjelaskan arsitektur parser generator.
  \item Saya dapat menulis deklarasi Bison.
  \item Saya dapat menulis aturan dan semantic action.
  \item Saya dapat membangun AST dari aksi.
  \item Saya dapat menangani konflik dengan precedence.
  \item Saya dapat mengintegrasikan Flex dan Bison.
  \item Saya dapat membandingkan Bison, ANTLR, dan parser tangan.
  \item Saya sudah menyelesaikan M3/Praktikum 3.
\end{checklist}

\section*{Rangkuman}
\addcontentsline{toc}{section}{Rangkuman}

Pada bab ini telah dibahas:
\begin{itemize}[leftmargin=*]
  \item arsitektur parser generator dan keuntungannya;
  \item deklarasi, aturan, dan konflik Bison;
  \item semantic action dan konstruksi AST;
  \item integrasi Flex--Bison dan build;
  \item pengujian parser, AST, dan pesan kesalahan sintaksis;
  \item perbandingan Bison dan parser tangan.
\end{itemize}

Bab XIV akan beralih ke analisis semantik: attribute grammar,
syntax-directed definition, dan translation scheme. AST yang
dihasilkan pada bab ini menjadi input untuk analisis tersebut.

\section*{Evaluasi}
\addcontentsline{toc}{section}{Evaluasi}

\begin{enumerate}[label=\arabic*.,leftmargin=*]
  \item Jelaskan alur kerja parser generator.
  \item Tuliskan deklarasi token dan tipe untuk ekspresi.
  \item Jelaskan peran \code{\%left} dan \code{\%right}.
  \item Tuliskan semantic action untuk aturan penjumlahan.
  \item Bagaimana AST dibangun dari aksi?
  \item Jelaskan integrasi \code{yylex} dan \code{yyparse}.
  \item Bandingkan Bison dan ANTLR dari sisi kontrol dan
        produktivitas.
  \item Mengapa laporan konflik harus dibaca?
\end{enumerate}

\subsection*{Kunci dan Rubrik Singkat}
\begin{itemize}[leftmargin=*]
  \item Soal 1--4: ketepatan deklarasi dan aksi (70\%),
        kelengkapan (30\%).
  \item Soal 5--6: pemahaman integrasi (60\%), contoh (40\%).
  \item Soal 7--8: analisis (60\%), argumentasi (40\%).
\end{itemize}


\ifSubfilesClassLoaded{%
  \bibliographystyle{apalike}%
  \bibliography{../../references}%
}{}

\end{document}
